\section{Finite demand graphs and reverse closure}
\label{sec:demand-graphs}

Real proofs are rarely simple chains.  A requested observable may depend on several
approximations, each produced from shared lower-level data, with proof obligations
that fan out differently from computation.  We therefore organize a proof as a fixed
acyclic graph.

\subsection{Graph data}

Let \(V=\{0,\ldots,m-1\}\) be nodes in topological order: every dependency of node
\(v\) has smaller index.  For each edge \(v\to u\), where \(u\) is a dependency of
\(v\), supply monotone dual demand transformers
\[
  r_{v,u},j_{v,u}:\N\to\N,
  \qquad r_{v,u}(d)\le j_{v,u}(d).
\]
A demanded node carries a pair \((R_v,J_v)\).  We adjoin a bottom value
\(\bot\), meaning ``this node is not used.''  This distinction is essential: numerical
demand \(0\) means that row zero is required, not that the node is unused.  Bottom is
below every pair; nonbottom pairs are ordered componentwise and joined by maximum.

Given a root \(z\) and requested output precision \(d\), initialize
\((R_z,J_z)=(d,d)\) and every other node to \(\bot\).  Traverse nodes in reverse
topological order.  Skip bottom nodes.  For each edge \(v\to u\) from a nonbottom node,
update
\[
\begin{split}
 R_u&\leftarrow\max\left(R_u,\max_{q\le R_v}r_{v,u}(q)\right),\\
 J_u&\leftarrow\max\left(J_u,\max_{q\le J_v}j_{v,u}(q)\right).
\end{split}
\]
Here joining a propagated pair into \(\bot\) returns that pair.  For multiple roots,
initialize each requested root and take joins.

\begin{theorem}[Reverse demand closure]
\label{thm:demand-closure}
The reverse pass computes the least bottom-extended demand assignment closed under
all edges leaving demanded nodes and dominating the root demand.  Its nonbottom nodes
are exactly the dependency cone of the root.  It is primitive recursive in the fixed
graph data and the final demand.  Every resulting pair satisfies \(R_v\le J_v\).
\end{theorem}

\begin{proof}
Because the order is topological, when node \(v\) is processed in reverse order, all
users of \(v\) have already propagated their final demands to it.  No later step can
increase \(v\)'s demand.  A node receives a nonbottom demand exactly when it lies on a
dependency path from the root.  Therefore one pass suffices and disconnected nodes do
not spuriously request row zero.

Closure follows directly from the update equations.  To prove minimality, let
\((R'_v,J'_v)\) be any closed assignment dominating the root.  Descend through the
reverse order.  Assume every already processed node has computed demand at most its
primed demand.  Monotonicity of prefix maxima and closure of the primed assignment show
that every propagation from the current node is at most the corresponding primed
dependency demand.  Taking maxima preserves the inequality.  Induction proves
\((R_v,J_v)\le(R'_v,J'_v)\) for all \(v\).

Primitive recursiveness follows because the graph is fixed and every update is a
bounded maximum of primitive-recursive functions.  If the graph itself is numerical
input, encode its finite adjacency table and run a bounded loop of length \(m\); the
same argument becomes a quantifier-free invariant about the table.

For \(R_v\le J_v\), induct on the reverse pass.  Initially the property holds.  On an
edge, \(R_v\le J_v\), the edge read transformer is bounded by its justification
transformer, and prefix maximum is monotone.  Hence the propagated read demand is at
most the propagated justification demand; componentwise joins preserve the property.
\end{proof}

\subsection{Why closure precedes induction}

The demand pass is run in the metatheory or as a preliminary primitive-recursive
computation, before the main correctness induction.  Its output is a finite vector of
bounds.  The proof then materializes exactly the table rows and contract instances
within those bounds.  This order avoids an illegal invariant saying ``for every future
precision, all needed rows are valid.''

In an object implementation, the demand pass itself is a bounded loop and may be
verified by quantifier-free induction.  The main trace uses its completed table as
fixed input.  Nested uses of quantifier-free induction are harmless: each invariant is
finite and explicit.

\subsection{Shared dependencies and proof reuse}

If two nodes depend on the same source, the closure takes the maximum of their demands.
The source table is computed once.  Likewise, its validity derivations may be stored
once in the proof schedule and cut into both consumers.  This is the proof-theoretic
analogue of common-subexpression elimination.

A graph representation is therefore not merely convenient notation.  It can reduce
proof terms and makes dependency auditing possible.  Each node declares exactly which
earlier formulas its local proof uses; undeclared global assumptions cannot enter
silently.

\subsection{Cycles and feedback}

A genuine cyclic dependency cannot be handled by a topological pass.  In practice one
of three situations occurs.

\begin{enumerate}[label=(\roman*),leftmargin=2.4em]
\item The apparent cycle is a bounded iteration.  Collapse it into one clocked loop
node whose state contains all mutually recursive values.
\item The cycle is a fixed-point specification with a contraction modulus.  Replace it
by a finite Picard iteration whose clock is derived from the modulus.
\item The cycle expresses an infinitary or impredicative construction with no supplied
majorant.  It lies outside the unconditional \(\FRNF\) interface and must remain a
boundary.
\end{enumerate}

Simultaneous primitive recursion, already available in the checked repository library,
is the standard implementation of the first case.

\subsection{Object encoding}

For a fixed small graph, the cleanest formalization generates one tuple of demand terms
per node and proves the closure equations directly.  For a graph supplied as finite
data, encode:

\begin{itemize}
\item the node count and topological order;
\item edge offsets in a flattened adjacency table;
\item codes selecting from a fixed finite family of demand transformers;
\item an active bit distinguishing bottom from the genuine numerical demand zero; and
\item the read/justification arrays updated by primitive recursion.
\end{itemize}

The interpreter is admissible because the family of node and edge operations is fixed
and finite; it is not a universal evaluator for arbitrary primitive-recursive syntax.
The object theorem \code{demandClosure_correct} in the Lean roadmap is the intended
closed derivation of this finite table algorithm.
